\chapter{What the bank would own and operate}
\label{ch:implementation}

The implementation has to answer a practical question: who can trust a decision,
and what can they inspect when the decision is challenged? A compliance
reviewer owns the meaning of a circular. A business owner sets the perimeter in
which that meaning will be used. Operations supplies and records the facts.
Engineering makes the approved meaning executable. A later incident may require
all four people to reconstruct one client or product decision.

This chapter follows that responsibility through one change. It starts with the
reviewer's view of the decision, then follows the evidence into a tested Java
control and a release. The detailed service paths, storage rules and task history
remain in Appendix~\ref{app:full-implementation-record} of the technical companion; they are reference
material for people who must build or audit the route.

\section{Begin with the decision a person needs to make}

A reviewer opening the SPI example should see the source passage, the client's
financial and category evidence, the alternatives that were considered, and the
case that would change the answer. The reviewer should not have to reconstruct
that explanation from a list of fields. If the portfolio is missing but the
alternative net-asset route is established, the screen should say why the missing
value does not block this particular conclusion. If two records disagree, it
should show both records and identify the person who must resolve them.

The same view serves different people in different ways. Compliance needs the
source and adopted reading. Operations needs the action, owner and deadline.
Engineering needs the precise values and histories used to execute the rule. An
auditor needs to replay the result later. One maintained decision record can
support all four views if the reader starts with the conclusion and can open the
reason behind it.

\section{Follow one change from source to release}

The route is easier to understand as a sequence of questions than as a software
architecture. What did the circular say? Which activity and legal entity are in
scope? Which reading did the bank adopt? What facts does that reading require?
Which cases would expose an error? Does the program return the same result as the
reviewed rule? Who approved this exact version, and when may the bank use it?

\begin{figure}[htbp]
\centering
\begin{tikzpicture}[node distance=7mm and 7mm,
 box/.style={draw=teal,fill=pale,rounded corners,align=center,text width=35mm,minimum height=15mm,font=\small},
 arr/.style={-{Latex[length=2mm]},thick,draw=navy}]
\node[box] (source) {Source passage and\\ business scope};
\node[box,right=of source] (meaning) {Reviewed meaning\\ and open questions};
\node[box,right=of meaning] (facts) {Facts, cases and\\ evidence owners};
\node[box,below=of facts] (program) {Reference rule and\\ generated Java};
\node[box,left=of program] (approval) {Independent checks\\ and approvals};
\node[box,left=of approval] (use) {Supervised use\\ and replay};
\draw[arr] (source)--(meaning); \draw[arr] (meaning)--(facts);\draw[arr] (facts)--(program);
\draw[arr] (program)--(approval); \draw[arr] (approval)--(use);
\draw[arr] (approval.north)--(meaning.south);
\end{tikzpicture}
\caption{A bank release carries a reviewed meaning and its evidence through
execution. Each step answers a different question.}
\label{fig:reader-implementation-route}
\end{figure}

The prototype now exercises this route for public documents and synthetic facts.
Consider the gift review. The reviewer begins with the selected paragraph and
source edition, sees alternative readings and missing coverage, and examines a
concrete distinguishing case where the fact meanings can be aligned. Separate
control components keep their answers and open questions when Java is generated.
If the source changes, affected work is reconsidered while the earlier decision
remains reproducible.

That journey includes different kinds of progress. Retrieving the FAQ supplies
context. Running a case exposes a program difference. An authorized decision
adopts a meaning for the bank. The first two are available in local development;
the third requires the institution's own authority and evidence. The demonstrated
route therefore supplies draft controls for review, with the open questions
carried through to delivery.

\section{Keep the release tied to its meaning}

A tested program is only useful when the bank can identify the rule it implements.
The release therefore carries the source edition, the adopted interpretation,
the factual definitions, the supported dates and activities, the test evidence
and the exact Java package together. Changing one of those parts creates a new
version that needs its own review. A content identifier helps detect accidental
mixing of versions; it does not make an administrator, source or interpretation
authoritative.

The same principle governs a correction. If a client withdraws consent on
Tuesday, a Monday decision must remain reproducible with the information known
on Monday. A Wednesday amendment must show which cases and clients change. The
bank can replace a defective software package while preserving the legal version
that applied to an earlier transaction. Historical replay is therefore part of
the control's meaning, not a convenience added after deployment.

\section{What the prototype has and has not demonstrated}

The engineering record contains checks for selected rules, edge cases,
deliberate mutations, release identity, withdrawal, amendment and archive
restoration. The interpretation route now also generates alternatives and
checks their source support. Earlier rounds included live development calls;
later integration exercises used those retained responses and controlled
inputs, with unresolved interpretations retained through delivery.

The distinction matters for the next investment decision. The mechanism can now
be exercised, but no comparative study has measured how often it misses a
requirement on unfamiliar material, how quickly a reviewer can use its reports,
or how it behaves within the bank's identity and data controls. Public sources
and synthetic identities remain the basis of the local demonstrations.

Those unanswered questions belong to the evaluation and operating decisions in
the last two chapters. Keeping them separate prevents a successful build from
being mistaken for a legal approval or a production control.

\section{The boundary before a bank pilot}

Before a supervised pilot, the bank would have to name the legal entities,
products, client categories, source families and decision points included in the
study. It would need an authorised legal owner, data owners for each fact,
independent engineering and compliance review, and a response for unknown or
conflicting evidence. The pilot should run beside the existing process so that
staff can compare decisions and record material disagreements. A case that cannot
be supported should stop or go to a named reviewer; it should not be converted
into a plausible default.

The implementation record makes the route concrete. The management decision is
whether the bank is willing to supply the authority, evidence and supervision
that the current public demonstration cannot supply by itself.
